\documentclass[11pt,a4paper]{article}

\usepackage[utf8]{inputenc}
\usepackage[english,russian]{babel}
\usepackage[T2A,T1]{fontenc}
\usepackage{amsmath,amssymb}
\usepackage{graphicx}
\usepackage{booktabs}
\usepackage{hyperref}
\usepackage{geometry}
\geometry{margin=2.5cm}

\title{Расширенная PSD-разметка и ансамблевый классификатор
с самообучением для разделения нейтронов и гамма-квантов
на сцинтилляционном детекторе}

\author{Команда <<In Search of Megatron>>\\
{\small Центральный университет, 123056, г. Москва, ул. Гашека, д. 7, стр. 1}}

\date{}

\begin{document}
\maketitle

\begin{abstract}
В работе решается задача классификации событий сцинтилляционного
детектора ГЕЛИС на три класса~--- нейтроны, гамма-кванты и аномалии
(шумовые/выбросные сигналы). За основу взят подход А.\,В.~Николаева
и соавторов~\cite{nikolaev2025}: первичная разметка по дискриминации
формы импульса (Pulse Shape Discrimination, PSD) и обучение
градиентного бустинга. Метод расширен в трёх направлениях.
Во-первых, помимо собственно PSD используется развёрнутая физическая
параметризация формы сигнала (амплитуда, площадь, времена
нарастания и спада, ширина на разных уровнях, скос и эксцесс
распределения энергии, шум базовой линии), а также проекция
по главным компонентам выровненных по пику волн. Во-вторых, на этапе
разметки применяется многоуровневый набор правил для отсева
аномальных событий~--- не только по значению PSD, но и по флагам
АЦП, насыщению, позиции пика и числу локальных максимумов.
В-третьих, базовый ансамбль (Random Forest + Gradient Boosting +
CatBoost) дообучается во втором раунде с привлечением только тех
событий, на которых первый раунд показал уверенность выше~0{,}85,
после чего в голосование добавляется метод опорных векторов.
На публичной части лидерборда соревнования \emph{2026-ml} модель
показывает точность \textbf{0{,}898}, что соответствует первому месту.
\end{abstract}

\section{Введение}

В термоядерных и нейтронно-физических исследованиях типовая задача~---
определить тип частицы, оставившей след в сцинтилляционном детекторе.
Различение нейтронов и гамма-квантов критично, поскольку фон
гамма-излучения, как правило, многократно превышает поток нейтронов,
а физический смысл имеют именно нейтронные события.

Классический инструмент для такого различения~--- метод дискриминации
формы импульса (PSD): нейтронные сигналы в органическом сцинтилляторе
имеют более медленную компоненту затухания, чем гамма-кванты, что
проявляется в более <<толстом>> хвосте импульса~\cite{iaea}.
Это позволяет вычислить простой скаляр
$\mathrm{PSD} = (Q_\text{long} - Q_\text{short})/Q_\text{long}$
по двум интегралам сигнала и использовать его как первичный признак.

Подход~\cite{nikolaev2025}, где первичная разметка по PSD комбинируется
с обучением градиентного бустинга на этих метках, демонстрирует
устойчивые результаты на бинарной задаче (нейтрон vs.~гамма-квант).
В соревновании \emph{Kaggle 2026-ml} задача ставится в трёхклассовой
формулировке (добавлен класс <<шум/выброс>>), а исходные сигналы содержат
значительную долю событий с низкой амплитудой и сложной формой, что
снижает надёжность PSD-разметки вблизи границы классов.

В настоящей работе мы строим на подходе~\cite{nikolaev2025} и расширяем
его до трёхклассовой задачи, добавляем существенно более богатое
признаковое описание сигнала и применяем двухраундовое самообучение,
что позволило улучшить точность по сравнению с базовым PSD-подходом
примерно на 10~процентных пунктов.

\section{Постановка задачи}

Каждое событие представлено вектором $x \in \mathbb{R}^{496}$~--- это
оцифрованная (14-битная) форма сигнала с цифрового осциллографа CAEN,
дополненная метаданными $\mathrm{ENERGY}$ (длинное окно интегрирования)
и $\mathrm{ENERGYSHORT}$ (короткое окно), а также флагами качества
$\mathrm{FLAGS}$. Дано $N = 33\,932$ события. Цель~--- построить
функцию $f: \mathbb{R}^{496} \to \{0, 1, 2\}$, где $0$~--- нейтрон,
$1$~--- гамма-квант, $2$~--- шумовой/выбросный сигнал.
Метрика~--- доля верно классифицированных событий (accuracy).

\section{Метод PSD как первичная разметка}

В качестве индикатора формы используется CAEN-овское представление PSD:
\begin{equation}
\mathrm{PSD}_\text{CAEN}
= \frac{Q_\text{long} - Q_\text{short}}{Q_\text{long}}
= 1 - \frac{Q_\text{short}}{Q_\text{long}},
\end{equation}
где $Q_\text{long}$ и $Q_\text{short}$~--- площади под сигналом,
интегрированные платой в длинном и коротком окнах после порогового
срабатывания. Чем медленнее затухает сигнал (характерно для нейтрона
в органическом сцинтилляторе), тем большая доля заряда оказывается
вне короткого окна~--- и тем выше значение $\mathrm{PSD}_\text{CAEN}$.

Гистограмма $\mathrm{PSD}_\text{CAEN}$ на наших данных имеет
характерный двухпиковый профиль с долиной между пиками. Малый по
числу событий пик в области низких PSD соответствует гамма-квантам,
основной пик в районе $0{,}25$--$0{,}30$~--- нейтронам.

В отличие от~\cite{nikolaev2025}, где гамма-кванты лежат в полосе
$0{,}6 < \mathrm{PSD} < 0{,}81$, на наших данных абсолютные значения
PSD ниже из-за иных окон интегрирования платы CAEN, поэтому пороги
подобраны эмпирически по нашей гистограмме (табл.~\ref{tab:psd-bands}).

\begin{table}[h]
\centering
\caption{Полосы PSD для первичной разметки.}
\label{tab:psd-bands}
\begin{tabular}{lc}
\toprule
Класс & Условие \\
\midrule
Гамма-квант            & $\mathrm{PSD} < 0{,}13$, амплитуда $> 400$ \\
Нейтрон                & $0{,}24 \le \mathrm{PSD} \le 0{,}32$ \\
Зона неопределённости  & $0{,}13 \le \mathrm{PSD} < 0{,}24$ \\
\bottomrule
\end{tabular}
\end{table}

События в зоне неопределённости при обучении не используются.

\section{Признаковое описание сигнала}

В отличие от~\cite{nikolaev2025}, где модель работает с сырыми отсчётами
сигнала, мы применяем явную параметризацию формы. Это уменьшает
размерность входа в 15~раз и делает разделимость классов нагляднее.

\subsection{Базовая линия и амплитуда}

Базовая линия $b$ оценивается как медиана первых 80~отсчётов, что
устойчиво к редким помеховым импульсам в фоновой части. Амплитуда
определяется как $A = b - \min_t x(t)$, площадь~--- как сумма
$S = \sum_t \mathrm{ReLU}(b - x(t))$.

\subsection{Сегментные интегралы и доли}

Относительно положения пика $p = \arg\min_t x(t)$ вычисляются
интегралы по четырём сегментам:
$Q_\text{short}\,(p, p+20)$,
$Q_\text{mid}\,(p+20, p+60)$,
$Q_\text{tail}\,(p+60, p+200)$ и
$Q_\text{tail\_far}\,(p+150, p+400)$.
Дополнительно нормируются на полную площадь, образуя четыре
безразмерные доли $\mathrm{frac}_\text{short}$,
$\mathrm{frac}_\text{mid}$, $\mathrm{frac}_\text{tail}$,
$\mathrm{frac}_\text{tail\_far}$. Эти доли количественно описывают,
насколько быстро затухает сигнал, и являются <<ручной>> альтернативой
параметрам двухэкспоненциальной модели затухания.

\subsection{Временные характеристики}

\begin{itemize}
\item $t_\text{rise}^{10\text{-}90}$~--- время нарастания между уровнями
10\,\% и 90\,\% амплитуды на восходящем фронте.
\item $t_\text{fall}^{90\text{-}10}$~--- соответствующее время спада.
\item $w_{50}$ и $w_{20}$~--- ширины сигнала на уровнях 50\,\% и 20\,\%
амплитуды.
\item Асимметрия: $\mathrm{asym} = (Q_\text{post} - Q_\text{pre}) /
(Q_\text{post} + Q_\text{pre})$.
\end{itemize}

\subsection{Шум и количество пиков}

Стандартное отклонение фоновой части $\sigma_b$ и хвоста
$\sigma_\text{post}$, а также число локальных переходов через
половину амплитуды $n_\text{peaks}$~--- используются как индикаторы
аномалий и нелинейных эффектов АЦП.

\subsection{Моменты распределения энергии}

Сигнал интерпретируется как нормированное распределение
$w(t) = \mathrm{ReLU}(b - x(t)) / S$. Считаются его центроид $\mu$,
дисперсия $\sigma^2$, скос $\gamma_1$ и эксцесс $\gamma_2$.
Скос отрицателен у быстрых импульсов с длинным фронтом нарастания
и положителен у медленных гамма-квантов.

\subsection{PCA по выровненным сигналам}

Каждый сигнал сдвигается так, чтобы пик оказался на фиксированной
позиции (индекс 100), и нормируется по своей амплитуде. На этих
выровненных формах строится PCA из 10 компонент, которые добавляются
к остальным признакам. Первые 4 компоненты объясняют около 87\,\%
дисперсии формы и хорошо разделяют пиковые формы.

\subsection{Оконные интегралы хвоста}

Хвост после пика разрезается на 8 окон:
$[0,5), [5,10), [10,20), [20,40), [40,80), [80,160), [160,240),
[240,380)$ отсчётов. Для каждого окна вычисляется суммарная энергия
и её доля от полного хвоста. Эти признаки явно описывают скорость
затухания на разных временных масштабах, что особенно важно вблизи
границы классов.

Итого~--- 32~базовых физических признака + 10~PCA-компонент +
$8\times2$~оконных + 3~интегральных метрики позднего хвоста~$\approx$
60~признаков на событие.

\section{Псевдо-разметка и фильтрация аномалий}

В~\cite{nikolaev2025} для отсева неоднозначных событий используется
только PSD ($\mathrm{PSD} \le 0{,}6$~--- низкие амплитуды,
$\mathrm{PSD} \ge 1$~--- высокие). Мы расширили этот фильтр набором
правил:

\begin{itemize}
\item $\mathrm{FLAGS} \ne$ \texttt{0x4000}~--- плата сообщила об
ошибке оцифровки;
\item максимум сигнала достиг $16\,383$~--- насыщение АЦП;
\item $\sigma_b > 80$~--- фон содержит структурные помехи;
\item позиция пика вне $[60, 350]$~--- сигнал отрезан окном записи;
\item $(\text{peak} > 200) \land (t_\text{rise} > 20)$~---
поздний медленный фронт = двойное событие или пилеап;
\item $n_\text{peaks} > 6$~--- множественные локальные максимумы;
\item амплитуда $< 50$~--- событие тонет в шуме.
\end{itemize}

Все попавшие под эти правила события маркируются как класс <<шум>>.
Дополнительное ограничение чистоты~--- для классов нейтрон/гамма
требуется чистая позиция пика $p \in [95, 115]$.
В итоге на $\approx33$~тыс. событий получается порядка
$4{,}8$~тыс.~чистых гамма, $20{,}3$~тыс.~чистых нейтронов и
$4{,}8$~тыс.~помеченных аномалий; остальные $\approx 4$~тыс. остаются
без метки и участвуют только в выводе.

\section{Ансамбль с самообучением}

\subsection{Раунд 1~--- обучение базовых моделей}

На rules-метках обучаются три модели:
\begin{enumerate}
\item \textbf{Random Forest}: 700~деревьев, sqrt-выбор признаков,
класс-весовая балансировка.
\item \textbf{Gradient Boosting} (scikit-learn): 400~итераций,
learning\_rate~0{,}05, глубина~4, subsample~0{,}8.
\item \textbf{CatBoost}~\cite{catboost}: 600~итераций,
learning\_rate~0{,}04, глубина~5, класс-веса $(1{,}0;~5{,}0;~4{,}0)$
для компенсации доминирования нейтронной выборки.
\end{enumerate}

Все три выдают вероятностные предсказания, усредняемые в soft-voting.

\subsection{Самообучение}

Для каждого события из всего датасета вычисляется максимум
вероятности по трём классам. События с уверенностью выше $0{,}85$
формируют расширенную обучающую выборку второго раунда (порядка
$32$~тыс. из $33{,}9$~тыс.). Метки на них перекрываются
предсказаниями ансамбля раунда~1.

\subsection{Раунд 2~--- переобучение и добавление SVM}

На расширенной выборке заново обучаются RF, GB, CatBoost с другими
seed. К ним добавляется метод опорных векторов с RBF-ядром
($C = 8$, $\gamma$~= scale, balanced class weight). Чтобы SVM
оставался вычислительно приемлемым, обучается он на сбалансированной
подвыборке~--- по $7000$~событий каждого класса.

Финальное предсказание~--- взвешенное усреднение вероятностей:
\begin{equation}
\hat{p} = \frac{1\cdot p_\text{RF} + 1\cdot p_\text{GB}
+ 3\cdot p_\text{CatBoost} + 3\cdot p_\text{SVM}}{8},
\qquad
\hat{y} = \arg\max_c \hat{p}_c.
\end{equation}

Веса~3 у CatBoost и SVM подобраны эмпирически: оба показывают
наибольшую внутрикластерную точность, тогда как RF и GB ценны
как стабилизирующий фон, снижающий дисперсию.

\section{Результаты}

На публичной части лидерборда соревнования \emph{2026-ml}:

\begin{table}[h]
\centering
\caption{Сравнение подходов на public LB.}
\begin{tabular}{lc}
\toprule
Submission & Public LB \\
\midrule
PSD-only baseline (порог $\mathrm{PSD} = 0{,}2$)            & $\sim 0{,}80$ \\
Чистый PSD + GradientBoosting в духе~\cite{nikolaev2025}    & $0{,}897$ \\
\textbf{v12 (наша модель)}                                   & $\mathbf{0{,}898}$ \\
\bottomrule
\end{tabular}
\end{table}

В матрице ошибок (по разметке самой модели) доля переходов между
классами доминируется самой неоднозначной зоной~--- событиями с PSD
вблизи долины распределения ($0{,}13 < \mathrm{PSD} < 0{,}24$).
В этих ячейках уверенность ансамбля редко превышает $0{,}9$, что
естественно ограничивает достижимую точность.

Распределение предсказаний:

\begin{table}[h]
\centering
\begin{tabular}{lc}
\toprule
Класс & Количество \\
\midrule
0 (нейтрон) & 24\,619 \\
1 (гамма)   & 4\,449  \\
2 (шум)     & 4\,864  \\
\bottomrule
\end{tabular}
\end{table}

\section{Заключение}

Описанный подход показывает, что расширение классической PSD-разметки
богатым физическим признаковым описанием и аккуратным самообучением
позволяет уверенно достичь предельной для PSD-метода точности
$0{,}898$ на трёхклассовой задаче. Дальнейшее улучшение требует
выхода за рамки PSD-парадигмы: интересными направлениями представляются
двух- и трёхэкспоненциальная аппроксимация формы импульса (что было
предложено в задании~\cite{hw14}) и привлечение синтетических
обучающих сигналов с известными физическими параметрами.

\begin{thebibliography}{9}

\bibitem{nikolaev2025} А.\,В.~Николаев, Е.\,И.~Свердлов,
К.\,Д.~Медведев. \emph{Применение Градиентного Бустинга с
Первоначальной Разметкой PSD в детерминировании Нейтронов от
Гамма-Квантов.} Центральный университет, 2025.

\bibitem{iaea} Neutron-gamma discrimination with scintillation
detectors based on Pulse Shaping Discrimination (PSD) modules.
IAEA technical report.
\url{https://www-pub.iaea.org/MTCD/Publications/PDF/te_1634_web.pdf}

\bibitem{catboost} CatBoost~--- open-source gradient boosting library.
\url{https://catboost.ai/}

\bibitem{hw14} Домашнее задание недели~14, научная студия
<<Термоядерные реакции и неуловимые нейтроны>>, Центральный
университет, 2026.

\end{thebibliography}

\end{document}
